\documentclass[11pt,a4paper,oneside,openany]{book}
\usepackage{ugm-modul}
\renewcommand{\modulheader}{Panduan Pandora}
\usepackage{hyperref}
\hypersetup{hidelinks,
  pdftitle={Dari Rekaman ke Dokumen Resmi: Panduan Pengguna Pandora}, pdfauthor={Tim Pengelola Pandora}}

\begin{document}

% =====================================================================
% SAMPUL
% =====================================================================
\begin{titlepage}
\thispagestyle{empty}
\begin{tikzpicture}[remember picture, overlay]
  \fill[ugmnavy] (current page.south west) rectangle (current page.north east);
  \node[anchor=north west] at ([xshift=1.6cm, yshift=-1.4cm]current page.north west)
    {\includegraphics[width=5.2cm]{figures/ugm-wordmark-white.png}};

  \begin{scope}[shift={(current page.center)}]
    \fill[ugmgold, opacity=0.14] (2.2,7.4) circle (1.7cm);
    \fill[ugmgold, opacity=0.25] (2.2,7.4) circle (0.95cm);
    \fill[ugmgold] (2.2,7.4) circle (0.42cm);
    \draw[white, opacity=0.35, line width=0.6pt]
      (2.2,7.4) -- (-3.2,10.6) (2.2,7.4) -- (6.4,10.7) (2.2,7.4) -- (-4.8,5.2)
      (2.2,7.4) -- (7.1,5.0) (2.2,7.4) -- (-1.0,3.4) (2.2,7.4) -- (5.0,2.9);
    \draw[white, opacity=0.18, line width=0.5pt, dashed]
      (-3.2,10.6) -- (-4.8,5.2) (6.4,10.7) -- (7.1,5.0) (-1.0,3.4) -- (5.0,2.9) (-4.8,5.2) -- (-1.0,3.4);
    \foreach \x/\y in {-3.2/10.6, 6.4/10.7, -4.8/5.2, 7.1/5.0, -1.0/3.4, 5.0/2.9}
      \fill[white] (\x,\y) circle (0.16cm);
    \foreach \x/\y in {-7.2/9.3, 9.0/8.2, 0.6/11.6, 8.4/1.6, -6.1/2.0, 3.7/12.0}
      \fill[white, opacity=0.5] (\x,\y) circle (0.07cm);
    \node[white, font=\small\sffamily, anchor=south] at (-3.2,10.85) {Unggah};
    \node[white, font=\small\sffamily, anchor=south] at (6.4,10.95) {Transkrip};
    \node[white, font=\small\sffamily, anchor=east] at (-4.95,5.2) {Rapi AI};
    \node[white, font=\small\sffamily, anchor=west] at (7.35,5.0) {Notulensi};
    \node[white, font=\small\sffamily, anchor=north] at (-1.0,3.15) {Dokumen};
    \node[white, font=\small\sffamily, anchor=west] at (5.25,2.9) {Riwayat};
    \node[white, font=\small\sffamily\bfseries, anchor=north] at (2.2,6.05) {Pandora};
  \end{scope}

  \fill[ugmgold] ([xshift=1.6cm, yshift=-14.8cm]current page.north west) rectangle ([xshift=3.6cm, yshift=-14.6cm]current page.north west);
  \node[anchor=north west, text width=16.5cm, align=left, text=white]
    at ([yshift=-15.4cm, xshift=1.6cm]current page.north west) {%
    {\bfseries\fontsize{36}{42}\selectfont Dari Rekaman\\ke Dokumen Resmi}\\[12pt]
    {\bfseries\fontsize{16}{22}\selectfont\color{ugmgold} Setiap kata dari rekamanmu, lebih mudah ditelusuri}\\[8pt]
    {\fontsize{12.5}{18}\selectfont Panduan pengguna transkripsi, notulensi, dan dokumen resmi}\\[30pt]
    {\fontsize{10.5}{14}\selectfont\color{white!70} Modul Pelatihan Pandora untuk pengguna di lingkungan UGM\\
     Tim Pengelola Pandora\\
     Edisi November 2026}};

  \node[anchor=south west, fill=ugmgold, text=ugmnavy, font=\small\bfseries, inner xsep=12pt, inner ysep=6pt]
    at ([yshift=2.9cm, xshift=1.6cm]current page.south west) {ugm.ac.id};
  \fill[ugmnavydeep] (current page.south west) rectangle ([yshift=1.3cm]current page.south east);
  \node[text=ugmgold, font=\small\bfseries] at ([yshift=0.65cm]current page.south) {Merakyat, Mandiri, Berkelanjutan};
\end{tikzpicture}
\clearpage
\end{titlepage}

% =====================================================================
% BAGIAN AWAL
% =====================================================================
\frontmatter
\tableofcontents
\listoffigures
\listoftables
\chapter*{Glosarium}
\addcontentsline{toc}{chapter}{Glosarium}

Istilah asing dicetak miring di dalam teks. Daftar berikut disusun menurut abjad.

\begin{description}[leftmargin=4.4cm, style=nextline, itemsep=6pt]
  \item[Halusinasi (\textit{hallucination})] Keluaran model yang terdengar meyakinkan tetapi tidak sesuai dengan isi rekaman, misalnya kalimat yang berulang tanpa dasar.
  \item[Kecerdasan artifisial (\textit{artificial intelligence}, AI)] Teknologi yang memungkinkan komputer menyusun, merangkum, atau mengenali pola dari data. Pandora memakainya untuk merapikan teks dan menyusun naskah.
  \item[Kepala dokumen] Bagian awal naskah dinas yang memuat tanggal, tempat, unit, dan keterangan lain yang diisi pengguna. Pandora tidak mengarang isi kolom ini.
  \item[Keluaran] Berkas hasil yang dibuat Pandora, misalnya transkrip, notulensi, atau naskah dalam format tertentu.
  \item[Model] Program pengenal suara yang mengubah audio menjadi teks. Makin besar model, makin teliti, tetapi prosesnya makin lama.
  \item[Naskah dinas] Dokumen resmi instansi dengan bagian dan bahasa baku, misalnya laporan kegiatan atau notulen rapat.
  \item[Notulensi (\textit{minutes})] Catatan resmi jalannya rapat yang memuat ringkasan, pembahasan, keputusan, dan tindak lanjut.
  \item[Rapikan transkrip] Fungsi untuk menyusun ulang transkrip agar lebih mudah dibaca. Transkrip asli tidak berubah.
  \item[Riwayat] Daftar folder hasil milik akun Anda yang dapat dibuka kembali.
  \item[Segmen] Potongan teks pendek dengan waktu awal dan akhir dalam rekaman.
  \item[Transkrip] Teks hasil pengubahan rekaman audio.
  \item[Transkripsi (\textit{transcription})] Proses mengubah rekaman suara menjadi teks.
  \item[Unggah] Tindakan memasukkan berkas audio dari komputer ke Pandora.
\end{description}


% =====================================================================
% ISI
% =====================================================================
\mainmatter
\chapter{Pendahuluan}

\section{Latar Belakang}
Rapat, kegiatan lapangan, dan wawancara menghasilkan rekaman yang sering hanya disimpan, lalu jarang dibuka lagi. Ketika notulen atau laporan harus disusun, rekaman itu harus didengarkan ulang dari awal. Proses ini memakan waktu, dan kesalahan kecil dalam mencatat angka atau keputusan mudah terjadi. Di sisi lain, banyak instansi kini mencari cara agar hasil rekaman bisa dipakai sebagai bahan kerja yang cepat dan dapat diperiksa ulang.

Pandora dibangun untuk menjawab kebutuhan itu. Aplikasi ini mengubah rekaman audio menjadi teks, lalu membantu menyusun notulensi dan naskah dinas dari teks tersebut. Penggunaannya berada dalam kerangka tata kelola layanan digital pemerintah (Presiden Republik Indonesia, 2018). Penggunaan kecerdasan artifisial juga perlu mengikuti etika yang ditetapkan kementerian terkait (Kementerian Komunikasi dan Informatika Republik Indonesia, 2023).

\section{Deskripsi Singkat}
Modul ini membahas Pandora dari langkah pertama sampai hasil akhir. Bab II menjelaskan cara masuk dan mengenali bagian-bagian layar. Bab III membahas cara mentranskripsi rekaman. Bab IV membahas cara membaca hasil dan merapikannya. Bab V membahas notulensi rapat, dan Bab VI membahas dokumen resmi. Bab VII membahas riwayat hasil dan masalah yang sering muncul. Penutup mengajak Anda kembali ke pertanyaan pemantik di bab ini.

Setiap bab materi dibuka dengan indikator keberhasilan, dilanjutkan uraian yang disertai gambar atau tabel, lalu ditutup dengan latihan dan rangkuman. Untuk mengikat seluruh pembahasan, modul ini memakai satu ilustrasi kasus yang muncul di setiap bab. Ilustrasi itu bersifat rekaan, tetapi disusun dari jenis pekerjaan yang lazim dijumpai di lingkungan unit kerja.

\begin{ilustrasi}
Rina adalah staf sekretariat sebuah unit kerja. Pada suatu hari ia menghadiri tiga rapat, dan kepala unit meminta notulen ketiganya sudah siap sebelum rapat koordinasi besok pagi. Ia belum pernah memakai Pandora, dan ia tidak ingin salah mencatat keputusan yang dibaca banyak orang.

Bab-bab berikutnya mengikuti langkah Rina, mulai dari masuk ke aplikasi sampai ia menyusun notulen dan dokumen pertamanya.
\end{ilustrasi}

\section{Tujuan Pembelajaran}
\subsection{Hasil Belajar}
Setelah mengikuti pelatihan ini, peserta mampu menghasilkan transkrip, notulensi, dan dokumen resmi dari rekaman melalui Pandora, serta memeriksa hasil itu sebelum dibagikan.

\subsection{Indikator Hasil Belajar}
Peserta mampu:
\begin{enumerate}
  \item mengenali bagian layar Pandora dan membaca status layanan (Bab II);
  \item mentranskripsi satu rekaman dengan model dan bahasa yang sesuai (Bab III);
  \item membedakan transkrip asli dari transkrip rapi dan memeriksa angka terhadap transkrip asli (Bab IV);
  \item menyusun notulensi dan memeriksa keputusan, nama, serta tenggat terhadap transkrip (Bab V);
  \item menyusun dokumen resmi dengan kepala yang diisi sendiri dan memahami alasan naskah ditolak pemeriksa (Bab VI);
  \item mengelola riwayat hasil dan mengatasi masalah umum secara mandiri (Bab VII).
\end{enumerate}

\section{Materi Pokok dan Alokasi Waktu}
Modul ini dirancang untuk pelatihan selama 4 jam 20 menit dalam dua sesi. Tabel~\ref{tab:materi} memuat materi pokok setiap bab beserta perkiraan alokasi waktunya.

\begin{table}[H]
\caption{Materi pokok dan alokasi waktu}
\label{tab:materi}
{\small
\begin{tabularx}{\linewidth}{@{}K{1.6cm}YP{3.4cm}@{}}
\toprule
\rowcolor{ugmpanel}\thd{Bab} & \thd{Materi pokok} & \thd{Alokasi waktu} \\
\midrule
Bab I & Pendahuluan dan pertanyaan pemantik & 20 menit \\
Bab II & Masuk dan mengenal Pandora & 30 menit \\
Bab III & Transkripsi rekaman & 45 menit \\
Bab IV & Membaca dan merapikan hasil & 40 menit \\
Bab V & Notulensi rapat & 45 menit \\
Bab VI & Dokumen resmi & 45 menit \\
Bab VII & Riwayat dan masalah umum & 20 menit \\
Penutup & Kembali ke pertanyaan pemantik & 15 menit \\
\bottomrule
\end{tabularx}
}
\end{table}

\section{Petunjuk Belajar}
Setiap bab dibuka dengan indikator keberhasilan, lalu uraian materi yang diselingi ilustrasi kasus, latihan, dan rangkuman. Istilah teknis yang belum Anda kenal dapat dicari pada Glosarium di bagian awal modul. Untuk latihan, Anda memerlukan akun UGM yang sudah aktif di Pandora dan berkas contoh audio pendek yang disediakan pengelola pelatihan.

Informasi tentang fitur dalam modul ini diperiksa terhadap aplikasi pada 16 November 2026. Fitur dan pengaturan Pandora dapat berubah setelah tanggal itu. Bila tampilan di layar berbeda dengan modul, ikuti tampilan yang tersedia di aplikasi dan beri tahu pengelola agar modul diperbarui.

\section{Ketentuan Penggunaan Data}
Selama pelatihan, gunakan hanya dokumen publik atau rekaman contoh yang disediakan. Rekaman rapat internal, data pribadi, dan dokumen berklasifikasi tidak boleh diunggah sebelum ada persetujuan dari pimpinan unit dan kejelasan bahwa layanan yang dipakai sudah sesuai ketentuan. Ketentuan ini dibahas lagi pada Bab II dan Bab III, ketika Anda mengunggah rekaman pertama.

\section{Pertanyaan Pemantik}
Sebelum masuk ke materi, tuliskan jawaban Anda atas pertanyaan berikut, paling banyak lima butir. Jawaban yang sama akan Anda tinjau kembali pada Penutup.

\begin{quote}
\itshape ``Dari rapat atau kegiatan Anda minggu ini, apa yang paling sulit dicatat ulang? Siapa yang paling membutuhkan catatan itu? Bagaimana Anda tahu catatan itu benar?''
\end{quote}

\chapter{Masuk dan Mengenal Pandora}

Bab ini membahas langkah pertama yang Anda lakukan setiap kali memakai Pandora. Anda akan masuk dengan akun UGM, membaca status layanan di bagian atas layar, dan mengenali kelompok panel yang akan Anda pakai di bab-bab berikutnya. Rina memulai justru dari sini, sebab ia tidak ingin menunggu di tengah rapat tanpa tahu apakah layanan siap dipakai.

\begin{indikator}
Setelah mempelajari bab ini, peserta mampu:
\begin{enumerate}
  \item masuk ke Pandora dengan akun UGM yang sah;
  \item membaca status layanan di bagian atas layar dan menentukan apakah fitur AI siap dipakai;
  \item menyebutkan fungsi setiap kelompok panel pada layar utama.
\end{enumerate}
\end{indikator}

\section{Masuk dengan Akun UGM}
Pandora menggunakan login Google. Akun yang diterima adalah akun dengan domain \texttt{ugm.ac.id} atau \texttt{mail.ugm.ac.id}. Setelah memilih akun, Anda diarahkan kembali ke halaman utama Pandora dan namanya tampil di layar.

Akun Anda menentukan kepemilikan hasil. Setiap transkrip, notulensi, dan dokumen yang Anda buat hanya tampil di riwayat akun itu. Karena itu, jangan membagikan komputer yang sedang terbuka di halaman Pandora kepada orang lain, dan keluar dari akun ketika selesai bekerja di komputer bersama.

\section{Membaca Status Layanan}
Bagian atas layar memperlihatkan status komponen yang dibutuhkan. Status ini membantu Anda tahu apa yang bisa dipakai sebelum memulai pekerjaan. Tabel~\ref{tab:status} merangkum arti setiap status.

\begin{table}[H]
\caption{Status layanan pada bagian atas layar}
\label{tab:status}
{\small
\begin{tabularx}{\linewidth}{@{}P{4.2cm}Y@{}}
\toprule
\rowcolor{ugmpanel}\thd{Komponen} & \thd{Arti dan tindakan} \\
\midrule
Pengolah audio & Komponen yang mengubah format rekaman. Bila bermasalah, transkripsi tidak bisa dimulai. Laporkan kepada pengelola. \\
Model transkripsi & Model pengenal suara tersedia. Bila tidak tersedia, pilih model lain yang tersedia atau minta pengelola memeriksanya. \\
Layanan AI & Diperlukan untuk merapikan transkrip, notulensi, dan dokumen resmi. Bila belum siap, tombol AI tidak dapat diklik. Transkripsi tetap bisa dipakai. \\
\bottomrule
\end{tabularx}
}
\end{table}

Tombol untuk merapikan transkrip, membuat notulensi, dan menyusun dokumen resmi hanya aktif bila layanan AI siap dan sebuah folder hasil sudah dipilih. Sebelum mengandalkan fitur AI untuk rapat yang penting, periksa status ini terlebih dahulu. Bila status tidak sesuai, catat waktunya dan beri tahu pengelola sebelum rapat berjalan.

\section{Mengenal Tata Letak Layar}
Layar utama Pandora terbagi menjadi beberapa kelompok panel. Gambar~\ref{fig:tata-letak} menunjukkan posisi kelompok utama, dan Tabel~\ref{tab:panel} menjelaskan fungsinya.

\begin{figure}[H]
\centering
\begin{tikzpicture}
  \node[kotakgelap, minimum width=15cm, minimum height=9mm] (h) at (0,0) {Bagian atas: status layanan};
  \node[kotak, minimum width=7.3cm, minimum height=24mm] (u) at (-3.9,-1.7) {Panel unggah dan pengaturan\\(berkas audio, model, bahasa)};
  \node[kotakemas, minimum width=7.3cm, minimum height=24mm] (s) at (3.9,-1.7) {Panel status, log, dan hasil\\(tab Asli sampai Dokumen)};
  \node[kotak, minimum width=15cm, minimum height=10mm] (r) at (0,-3.6) {Panel riwayat (folder hasil milik akun Anda)};
  \draw[garis] (h) -- (u);
  \draw[garis] (h) -- (s);
\end{tikzpicture}
\caption{Tata letak kelompok panel pada layar utama Pandora}
\label{fig:tata-letak}
\end{figure}

Panel yang ditandai emas adalah tempat Anda membaca hasil, sedangkan panel unggah menjadi titik awal setiap pekerjaan.

\begin{table}[H]
\caption{Fungsi setiap panel}
\label{tab:panel}
{\small
\begin{tabularx}{\linewidth}{@{}P{4.2cm}Y@{}}
\toprule
\rowcolor{ugmpanel}\thd{Panel} & \thd{Fungsi} \\
\midrule
File audio & Memilih berkas rekaman dengan cara diseret atau diklik. \\
Pengaturan & Memilih model dan bahasa. Pilihan terakhir tersimpan. \\
Status dan log & Menampilkan progres, pesan, dan potongan teks secara langsung. \\
Hasil & Menampilkan tab Asli, JSON, Rapi AI, Notulensi, dan Dokumen. \\
Riwayat & Membuka kembali hasil lama dan menghapus hasil yang tidak diperlukan. \\
\bottomrule
\end{tabularx}
}
\end{table}

\begin{ilustrasi}
Rina membuka Pandora pada pagi hari dan melihat status layanan AI belum siap. Tombol untuk membuat notulensi berwarna pudar. Ia tidak menunda pekerjaan. Ia tetap menyiapkan rekaman pertama untuk ditranskripsi, lalu mencatat jam saat status itu ia lihat, supaya bisa ia laporkan kepada pengelola bila masih belum siap di siang hari.
\end{ilustrasi}

\begin{kegiatan}{Latihan 1. Masuk dan membaca layar (10 menit)}
Masuk dengan akun UGM Anda. Tulis status tiga komponen di bagian atas layar, lalu tentukan apakah tombol AI dapat dipakai hari ini. Jika tombol AI tidak aktif, catat penyebab yang paling mungkin menurut Tabel~\ref{tab:status}. Keluaran yang diharapkan: satu kalimat yang menyatakan kesiapan layanan dan satu kalimat tindakan yang akan Anda ambil.
\end{kegiatan}

\begin{rangkuman}
Pandora dibuka dengan akun UGM, dan hasil Anda hanya tampil di riwayat akun itu. Status layanan di bagian atas layar menentukan fitur mana yang siap dipakai, dan tombol AI tidak aktif bila layanan AI belum siap. Layar terdiri dari panel unggah, status dan hasil, serta riwayat. Pada bab berikutnya, Rina mulai dengan rekaman rapat pertamanya.
\end{rangkuman}

\chapter{Transkripsi Rekaman}

Transkripsi adalah langkah yang mengubah rekaman menjadi teks. Hasilnya menjadi dasar bagi semua fitur lain di Pandora, sehingga kualitas transkripsi menentukan kualitas notulensi dan dokumen yang dibuat kemudian. Bab ini membahas cara memilih berkas, mengatur model dan bahasa, memulai dan memantau proses, serta membatalkannya bila perlu. Rina menerapkan langkah ini pada rekaman rapat pertamanya yang berdurasi 90 menit.

\begin{indikator}
Setelah mempelajari bab ini, peserta mampu:
\begin{enumerate}
  \item memilih berkas audio dengan format yang didukung dan menyebutkan alasan memilih model tertentu;
  \item memulai transkripsi dan membaca progres pada panel status serta log;
  \item membatalkan transkripsi dan menjelaskan apa yang tersisa setelah pembatalan.
\end{enumerate}
\end{indikator}

\section{Memilih Berkas Audio}
Pandora menerima berkas dengan format \texttt{.mp3}, \texttt{.wav}, \texttt{.m4a}, \texttt{.flac}, \texttt{.aac}, dan \texttt{.ogg}. Anda dapat menyeret berkas ke kotak unggah atau mengkliknya untuk memilih dari komputer. Setelah berkas terpilih, nama dan ukurannya tampil di bawah kotak.

Rekaman panjang dapat diproses. Pandora dirancang untuk rekaman rapat hingga sekitar tiga jam. Rekaman yang lebih panjang tetap dapat diproses, tetapi waktunya bertambah dan Anda perlu menyediakan waktu lebih banyak sebelum batas rapat berikutnya.

Bila Anda salah memilih berkas, klik \emph{Hapus file}. Tindakan ini hanya menghapus pilihan di halaman Anda. Berkas di komputer dan hasil di server tidak ikut terhapus.

\section{Memilih Model dan Bahasa}
Pengaturan transkripsi memuat dua pilihan. Tabel~\ref{tab:model} menjelaskan pilihan model.

\begin{table}[H]
\caption{Pilihan model transkripsi}
\label{tab:model}
{\small
\begin{tabularx}{\linewidth}{@{}K{2.4cm}YP{4.2cm}@{}}
\toprule
\rowcolor{ugmpanel}\thd{Model} & \thd{Ciri} & \thd{Saran pemakaian} \\
\midrule
\texttt{tiny} & Paling cepat, ketelitian paling rendah & Uji coba singkat \\
\texttt{base} & Cepat, ketelitian terbatas & Rekaman pendek dengan suara jelas \\
\texttt{small} & Seimbang antara kecepatan dan ketelitian & Rapat biasa \\
\texttt{medium} & Paling teliti, paling lama & Rapat penting, istilah teknis, atau banyak nama \\
\bottomrule
\end{tabularx}
}
\end{table}

Model transkripsi yang dipakai Pandora berakar pada pengenal suara berskala besar yang dilatih dengan data yang beragam (Radford dkk., 2023). Karena itu, ketelitiannya tetap bergantung pada kualitas rekaman. Bahasa yang tersedia adalah Indonesia (\texttt{id}) dan Inggris (\texttt{en}). Pilih bahasa sesuai bahasa yang paling banyak dipakai dalam rekaman. Pilihan terakhir Anda tersimpan dan dipakai lagi pada kunjungan berikutnya, sehingga Anda tidak perlu mengatur ulang setiap kali.

\section{Memulai dan Memantau Transkripsi}
Langkah kerjanya sebagai berikut.
\begin{enumerate}
  \item Pilih berkas audio.
  \item Periksa model dan bahasa pada panel pengaturan.
  \item Klik \emph{Transcribe}. Berkas diunggah ke server, lalu transkripsi berjalan di latar belakang.
  \item Pantau panel status. Persentase dan pesan progres tampil di sana.
  \item Lihat panel log. Potongan teks muncul per segmen, sehingga Anda bisa memeriksa apakah suara sudah dikenali dengan baik.
\end{enumerate}

Pandora mengerjakan satu transkripsi pada satu waktu untuk seluruh pengguna. Bila Anda mencoba memulai transkripsi saat pengguna lain sedang memproses rekaman, Pandora menampilkan pesan bahwa worker sedang digunakan. Tunggu sampai pekerjaan itu selesai, lalu coba lagi.

\section{Membatalkan Transkripsi}
Klik \emph{Batal} bila hasil sudah jelas keliru, misalnya berkas yang salah dipilih, atau bila proses terlalu lama. Jangan mengandalkan hasil dari pekerjaan yang dibatalkan. Periksa dulu folder hasilnya di riwayat, lalu mulai ulang dari langkah pemilihan berkas bila Anda masih membutuhkan transkripsinya.

\begin{ilustrasi}
Rina mengunggah rekaman rapat berdurasi 90 menit dengan model \texttt{small} dan bahasa Indonesia. Di panel log ia melihat beberapa bagian tidak menghasilkan teks, dan ia sempat mengira mesin berhenti. Ia memeriksa rekaman pada menit yang sama dan menemukan bahwa di sana memang ada jeda hening, ketika peserta mengambil dokumen. Hasil transkripsi kemudian dilanjutkan tanpa kalimat yang berulang. Dari pengalaman itu, ia belajar bahwa log membantu, tetapi keputusan tetap harus dicocokkan dengan rekaman asli.
\end{ilustrasi}

\begin{kegiatan}{Latihan 2. Mentranskripsi berkas contoh (15 menit)}
Unggah berkas contoh \texttt{sample\_75s.mp3} dengan bahasa Indonesia dan model \texttt{base}. Catat berapa lama proses selesai, berapa potongan teks yang muncul di log, dan apakah ada bagian yang berulang. Ulangi dengan model \texttt{small}, lalu bandingkan waktu dan ketelitian kata yang Anda catat. Keluaran yang diharapkan: tabel dua baris berisi model, waktu proses, dan satu catatan kualitas.
\end{kegiatan}

\begin{rangkuman}
Transkripsi dimulai dengan memilih berkas audio dalam format yang didukung, lalu mengatur model dan bahasa. Panel status dan log membantu Anda memantau proses dan memeriksa potongan teks. Hanya satu transkripsi yang berjalan pada satu waktu, dan pembatalan tidak menghasilkan transkrip lengkap. Rina tidak menyerah pada log yang tampak kosong. Ia memeriksa rekamannya sendiri, lalu melanjutkan ke tahap merapikan hasil.
\end{rangkuman}

\chapter{Membaca dan Merapikan Hasil}

Setelah transkripsi selesai, Pandora menyediakan beberapa bentuk hasil. Tidak semuanya sama. Transkrip asli adalah sumber, sedangkan transkrip rapi, notulensi, dan dokumen resmi adalah turunan darinya. Bab ini membahas cara membaca setiap tab, cara merapikan transkrip dengan aman, dan cara memeriksa hasil turunan terhadap sumbernya. Pemeriksaan ini penting, sebab setiap turunan bisa menyimpang dari rekaman bila tidak diperiksa.

\begin{indikator}
Setelah mempelajari bab ini, peserta mampu:
\begin{enumerate}
  \item membedakan fungsi tab Asli, JSON, dan Rapi AI;
  \item merapikan transkrip dengan fungsi \emph{Rapikan transkrip} tanpa mengubah transkrip asli;
  \item memeriksa angka dan kutipan pada hasil turunan terhadap transkrip asli.
\end{enumerate}
\end{indikator}

\section{Membaca Tab Hasil}
Panel hasil memuat beberapa tab. Tabel~\ref{tab:tab} menjelaskan fungsi setiap tab dan sumbernya.

\begin{table}[H]
\caption{Tab pada panel hasil}
\label{tab:tab}
{\small
\begin{tabularx}{\linewidth}{@{}P{3.2cm}YP{3.8cm}@{}}
\toprule
\rowcolor{ugmpanel}\thd{Tab} & \thd{Isi dan fungsi} & \thd{Sumber} \\
\midrule
Asli & Transkrip mentah, sesuai hasil transkripsi & Rekaman \\
JSON & Segmen dan metadata untuk keperluan teknis, termasuk waktu awal dan akhir setiap potongan & Rekaman \\
Rapi AI & Transkrip yang dirapikan agar lebih mudah dibaca & Transkrip asli \\
Notulensi & Notulensi rapat yang tersusun dengan bagian ringkasan, pembahasan, keputusan, dan tindak lanjut & Transkrip asli \\
Dokumen & Naskah dinas resmi dengan kepala yang diisi pengguna & Transkrip asli \\
\bottomrule
\end{tabularx}
}
\end{table}

Tab Asli adalah rujukan utama. Bila ada keraguan, kembali ke tab ini. Tab JSON berguna bila Anda perlu melihat di menit berapa sebuah kalimat diucapkan, sebab setiap segmen menyimpan waktunya.

Gambar~\ref{fig:turunan} menunjukkan hubungan antara transkrip asli dan hasil turunannya.

\begin{figure}[H]
\centering
\begin{tikzpicture}
  \node[kotakgelap, minimum width=4.2cm, minimum height=16mm] (r) at (0,0) {Rekaman audio};
  \node[kotakgelap, minimum width=4.2cm, minimum height=16mm] (a) at (5.2,0) {Transkrip asli\\(tab Asli dan JSON)};
  \node[kotak, minimum width=4.2cm, minimum height=13mm] (b) at (10.4,2.1) {Transkrip rapi AI};
  \node[kotak, minimum width=4.2cm, minimum height=13mm] (c) at (10.4,0) {Notulensi};
  \node[kotakemas, minimum width=4.2cm, minimum height=13mm] (d) at (10.4,-2.1) {Dokumen resmi};
  \draw[garis, ->] (r) -- (a);
  \draw[garis, ->] (a) -- (b);
  \draw[garis, ->] (a) -- (c);
  \draw[garis, ->] (a) -- (d);
\end{tikzpicture}
\caption{Hubungan transkrip asli dengan hasil turunannya}
\label{fig:turunan}
\end{figure}

Pada gambar, transkrip asli menjadi sumber bagi tiga turunan. Dokumen resmi ditandai emas karena ia paling sering dibagikan ke pihak lain. Turunan tidak boleh lebih dipercaya daripada sumbernya.

\section{Merapikan Transkrip}
Fungsi \emph{Rapikan transkrip} membuat versi transkrip yang lebih rapi. Langkahnya sebagai berikut.
\begin{enumerate}
  \item Buka hasil yang ingin Anda rapikan, lalu klik \emph{Rapikan transkrip}.
  \item Tunggu status selesai.
  \item Buka tab Rapi AI untuk membaca hasilnya.
\end{enumerate}

Fungsi ini tidak mengubah transkrip asli. Hasil rapi disimpan sebagai berkas terpisah, sehingga Anda dapat selalu kembali ke tab Asli. Hasil rapi juga dapat diunduh dalam format Word.

Walaupun lebih enak dibaca, hasil rapi tetap perlu diperiksa. Perapian yang baik menjaga makna, tetapi perapian yang kurang hati-hati dapat menghapus kalimat penting atau mengubah angka. Karena itu, untuk keperluan resmi, selalu rujuk transkrip asli.

\section{Memeriksa Hasil Turunan}
Setiap turunan perlu diperiksa dengan tiga pertanyaan sederhana.
\begin{enumerate}
  \item Apakah angka yang tertulis di turunan juga ada di transkrip asli, dan disebut oleh siapa?
  \item Apakah kutipan atau pernyataan sama dengan yang diucapkan, tidak lebih kuat dan tidak lebih lemah?
  \item Apakah kalimat yang berupa usulan tidak berubah menjadi keputusan?
\end{enumerate}

Pemeriksaan ini sebaiknya dilakukan dengan mencari kata kunci di tab Asli. Bila ada angka atau nama, cari kemunculannya di transkrip, lalu baca kalimat di sekitarnya. Dengan cara ini, Anda memeriksa sumbernya, bukan hanya membaca hasil turunan.

\begin{ilustrasi}
Rina merapikan transkrip rapatnya dan membaca hasil rapi itu dengan lancar. Ketika membandingkan angka anggaran, ia menemukan bahwa angka pada hasil rapi tidak sama dengan angka yang tercantum di tab Asli. Ia mencari angka itu di transkrip asli dan mendapati penyaji menyebut dua angka berbeda untuk dua kegiatan. Rina menulis catatan perbedaan itu, lalu memakai transkrip asli sebagai rujukan.
\end{ilustrasi}

\begin{kegiatan}{Latihan 3. Merapikan dan memeriksa (15 menit)}
Dari hasil Latihan 2, klik \emph{Rapikan transkrip}. Pilih tiga kalimat yang paling penting dari tab Asli, lalu cari pasangannya di tab Rapi AI. Tulis apakah makna setiap pasangan sama. Jika ada perbedaan, tandai bagian mana yang berubah dan apakah perubahan itu memengaruhi keputusan. Keluaran yang diharapkan: tiga pasangan kalimat dan satu kesimpulan tentang kelayakan hasil rapi untuk dibagikan.
\end{kegiatan}

\begin{rangkuman}
Transkrip asli adalah sumber bagi semua hasil lain. Tab Asli dan JSON membantu Anda memeriksa apa yang benar-benar diucapkan, sedangkan Rapi AI disimpan terpisah dan tidak mengubah transkrip asli. Setiap angka dan keputusan pada turunan perlu dicocokkan dengan sumbernya. Rina menemukan perbedaan angka dan memakai transkrip asli sebagai rujukan. Pada bab berikutnya, ia menyusun notulensi dari hasil yang sudah diperiksa.
\end{rangkuman}

\chapter{Notulensi Rapat}

Notulensi adalah catatan resmi jalannya rapat. Pandora dapat membuat rancangan notulensi dari transkrip dengan fungsi \emph{Buat notulensi rapat}. Rancangan itu membantu Anda memulai, tetapi notulensi yang dibagikan tetap menjadi tanggung jawab notulis. Bab ini membahas isi rancangan, cara memeriksanya, dan cara mengunduh hasil dalam format Word.

\begin{indikator}
Setelah mempelajari bab ini, peserta mampu:
\begin{enumerate}
  \item membuat rancangan notulensi dari transkrip yang sudah diperiksa;
  \item memeriksa keputusan, nama, jabatan, dan tenggat pada rancangan terhadap transkrip asli;
  \item mengunduh notulensi dalam format Word dan menyimpannya di folder kerja Anda.
\end{enumerate}
\end{indikator}

\section{Membuat Rancangan Notulensi}
Langkah membuat rancangan sebagai berikut.
\begin{enumerate}
  \item Pastikan layanan AI siap pada bagian atas layar.
  \item Pilih folder hasil dari transkripsi yang ingin dibuatkan notulensinya.
  \item Klik \emph{Buat notulensi rapat}, lalu tunggu. Proses ini bisa memakan beberapa menit, terutama untuk rekaman panjang.
  \item Buka tab Notulensi untuk membaca pratinjau.
\end{enumerate}

Rancangan notulensi memuat bagian ringkasan, pembahasan, keputusan, dan tindak lanjut. Setiap bagian disusun dari isi transkrip, bukan dari pengetahuan umum, sehingga kalimat yang tidak ada dalam rekaman seharusnya tidak muncul. Meski demikian, pemeriksaan tetap wajib.

\section{Memeriksa Rancangan}
Pemeriksaan notulensi berfokus pada empat hal yang paling sering menimbulkan masalah, yaitu keputusan, nama, tenggat, dan angka. Tabel~\ref{tab:periksa} menjadi daftar periksa Anda.

\begin{table}[H]
\caption{Daftar periksa rancangan notulensi}
\label{tab:periksa}
{\small
\begin{tabularx}{\linewidth}{@{}P{3.6cm}Y@{}}
\toprule
\rowcolor{ugmpanel}\thd{Aspek} & \thd{Yang diperiksa terhadap transkrip asli} \\
\midrule
Keputusan & Apakah setiap keputusan benar-benar diambil dalam rapat? Usulan, pendapat, dan rencanya tidak boleh ditulis sebagai keputusan. \\
Nama dan jabatan & Apakah setiap nama dan jabatan disebut dalam rekaman? Jangan menambah peserta yang tidak hadir. \\
Tenggat dan penanggung jawab & Apakah tenggat dan penanggung jawab disebut? Bila tidak disebut, tulis bahwa belum ditetapkan. \\
Angka dan tanggal & Apakah setiap angka dan tanggal ada di transkrip, dan siapa yang menyebutnya? \\
\bottomrule
\end{tabularx}
}
\end{table}

Bila sebuah bagian tidak jelas dalam rekaman, hasilnya biasanya ditandai dengan \texttt{[tidak jelas]} atau sejenisnya. Jangan menebak isi bagian itu. Dengarkan kembali rekaman di menit yang sesuai, lalu perbaiki notulensi secara manual.

\section{Mengedit dan Mengunduh}
Rancangan dari Pandora dapat diedit di luar aplikasi. Unduh notulensi dalam format Word melalui tautan unduh yang muncul setelah rancangan selesai, lalu buka berkas itu di Microsoft Word atau aplikasi sejenis. Nama berkas mengikuti pola \texttt{NotulenAI} diikuti nama folder hasil.

Sebelum dibagikan, lakukan dua hal. Pertama, bandingkan pembuka dan penutup notulensi dengan rekaman, sebab bagian ini paling sering diedit. Kedua, tambahkan keterangan bahwa notulensi disusun dengan bantuan AI dari rekaman dan belum disahkan, bila instansi Anda memerlukannya.

\begin{ilustrasi}
Rina membuat rancangan notulensi untuk rapat pertama. Pada bagian keputusan, ia menemukan kalimat bahwa anggaran ``disetujui'' untuk kegiatan kedua. Ketika ia kembali ke transkrip, ternyata yang disetujui hanya kajian awal, sedangkan anggaran penuh masih menunggu rapat berikutnya. Rina mengganti kalimat itu dengan rumusan yang sesuai rekaman, lalu mencatat sumber menitnya di pinggir naskah.
\end{ilustrasi}

\begin{kegiatan}{Latihan 4. Membuat dan memeriksa notulensi (20 menit)}
Dari transkrip latihan sebelumnya, buat rancangan notulensi. Pilih tiga keputusan atau tindak lanjut dari rancangan, lalu cari masing-masing di transkrip asli. Untuk setiap item, tulis apakah ia keputusan, usulan, atau pendapat, dan apakah namanya disebut dalam rekaman. Keluaran yang diharapkan: tabel tiga baris dengan kolom butir, kategori, dan hasil pemeriksaan.
\end{kegiatan}

\begin{rangkuman}
Pandora membuat rancangan notulensi dari transkrip yang sudah diperiksa. Rancangan itu tetap harus dibaca terhadap rekaman, terutama untuk keputusan, nama, tenggat, dan angka. Usulan tidak boleh berubah menjadi keputusan, dan bagian yang tidak jelas tidak boleh ditebak. Setelah diperiksa, notulensi dapat diunduh dalam format Word. Rina memperbaiki satu keputusan yang keliru sebelum notulensinya dibagikan.
\end{rangkuman}

\chapter{Dokumen Resmi}

Dokumen resmi adalah naskah dinas yang disusun dari transkrip dengan bagian dan bahasa baku instansi. Pandora menyediakan dua jenis, yaitu Laporan Kegiatan dan Notulen Rapat. Berbeda dengan notulensi yang berfungsi sebagai bahan kerja, dokumen resmi biasanya diarsipkan atau diajukan kepada pejabat. Karena itu, aturan penyusunannya lebih ketat. Pandora menolak naskah yang tidak lolos pemeriksaan, dan Anda perlu memahami alasannya.

\begin{indikator}
Setelah mempelajari bab ini, peserta mampu:
\begin{enumerate}
  \item memilih jenis dokumen dan mengisi kepala dokumen tanpa mengarang isinya;
  \item menjelaskan bagian yang tersusun pada setiap jenis dokumen;
  \item memahami dan menindaklanjuti alasan naskah ditolak pemeriksa.
\end{enumerate}
\end{indikator}

\section{Memilih Jenis dan Mengisi Kepala}
Klik \emph{Susun dokumen resmi} pada tab Dokumen. Jendela isian muncul dengan tiga bagian: jenis dokumen, judul, dan kepala dokumen. Judul awal diambil dari nama folder hasil, dan Anda dapat mengubahnya. Kolom kepala berubah mengikuti jenis dokumen yang dipilih.

Tabel~\ref{tab:kepala} memuat kolom kepala untuk setiap jenis.

\begin{table}[H]
\caption{Kolom kepala setiap jenis dokumen}
\label{tab:kepala}
{\small
\begin{tabularx}{\linewidth}{@{}P{3.6cm}Y@{}}
\toprule
\rowcolor{ugmpanel}\thd{Jenis dokumen} & \thd{Kolom kepala} \\
\midrule
Laporan Kegiatan & Hari dan tanggal, tempat, unit penyusun, disampaikan kepada \\
Notulen Rapat & Hari dan tanggal, waktu, tempat, pimpinan rapat, notulis \\
\bottomrule
\end{tabularx}
}
\end{table}

Isi kepala dengan informasi yang Anda ketahui. Pandora tidak mengarang tanggal, tempat, atau nama pejabat. Kolom yang Anda kosongkan tetap kosong pada naskah. Jangan menambahkan jabatan atau nama yang tidak Anda yakini, sebab kepala dokumen dibaca pertama kali oleh penerima.

\section{Bagian yang Tersusun}
Setiap jenis dokumen memiliki urutan bagian yang tetap. Tabel~\ref{tab:bagian} menjelaskan isi setiap bagian.

\begin{table}[H]
\caption{Bagian pada setiap jenis dokumen}
\label{tab:bagian}
{\small
\begin{tabularx}{\linewidth}{@{}P{3.6cm}Y@{}}
\toprule
\rowcolor{ugmpanel}\thd{Jenis} & \thd{Bagian, berurutan} \\
\midrule
Laporan Kegiatan & I Dasar dan maksud pelaksanaan. II Pelaksanaan. III Hasil yang dilaporkan. IV Kendala dan catatan perbedaan. V Tindak lanjut. VI Penutup. \\
Notulen Rapat & I Ringkasan rapat. II Peserta dan pihak yang disebut. III Pokok pembahasan. IV Keputusan. V Tindak lanjut. VI Penutup. \\
\bottomrule
\end{tabularx}
}
\end{table}

Beberapa aturan penyusunan berlaku untuk seluruh naskah. Setiap angka wajib disertai keterangan siapa yang menyebutnya, dan angka yang tidak disebut dalam rekaman tidak boleh ditulis. Bagian yang tidak terdengar jelas ditulis sebagai \texttt{[tidak jelas]}. Usulan, pendapat, dan keputusan dibedakan dengan tegas. Penutup Laporan Kegiatan menyatakan bahwa dokumen memerlukan pengesahan pejabat berwenang.

Gambar~\ref{fig:alur-dokumen} menunjukkan langkah penyusunan dari pilihan jenis sampai keluaran.

\begin{figure}[H]
\centering
\begin{tikzpicture}
  \node[kotak, text width=2.2cm, minimum height=16mm] (a) at (0,0) {Pilih jenis dan judul};
  \node[kotak, text width=2.2cm, minimum height=16mm] (b) at (3.0,0) {Isi kepala};
  \node[kotak, text width=2.2cm, minimum height=16mm] (c) at (6.0,0) {Susun bagian berurutan};
  \node[kotakemas, text width=2.2cm, minimum height=16mm] (d) at (9.0,0) {Periksa cakupan dan bentuk};
  \node[kotak, text width=2.2cm, minimum height=16mm] (e) at (12.0,0) {Keluaran: MD, PDF, DOCX};
  \draw[garis, ->] (a) -- (b);
  \draw[garis, ->] (b) -- (c);
  \draw[garis, ->] (c) -- (d);
  \draw[garis, ->] (d) -- (e);
\end{tikzpicture}
\caption{Alur penyusunan dokumen resmi}
\label{fig:alur-dokumen}
\end{figure}

Langkah yang ditandai emas adalah pemeriksaan. Di sinilah Pandora memutuskan apakah naskah boleh dinyatakan selesai.

\section{Memahami Naskah yang Ditolak}
Pemeriksaan naskah memeriksa tiga hal. Pertama, cakupan sumber, yaitu apakah naskah memakai bagian rekaman yang memadai. Kedua, bentuk naskah, yaitu apakah urutan bagian, judul, dan tata penulisan sesuai aturan. Ketiga, keterlacakan, yaitu apakah naskah dapat dikaitkan kembali dengan transkrip yang menjadi sumbernya. Bila salah satu tidak terpenuhi, Pandora menampilkan alasan dan tidak menyatakan dokumen selesai.

Cara menanggapinya sebagai berikut. Baca alasan yang ditampilkan. Bila masalahnya ada pada bentuk, misalnya bagian yang terlewat, coba susun ulang. Bila masalahnya ada pada cakupan, periksa apakah transkrip terpotong atau tidak lengkap, lalu ulangi transkripsi bila perlu. Jangan menyunting naskah yang ditolak lalu menyatakannya sah tanpa pemeriksaan ulang.

\section{Hasil Dokumen}
Setelah berhasil, Pandora menyediakan tiga berkas. Berkas \texttt{.md} berisi teks mentah. Berkas \texttt{.pdf} siap dicetak, tetapi hanya tersedia bila perkakas pembuat PDF terpasang di server. Berkas \texttt{.docx} dapat diedit di Word, sehingga cocok bila ada penyuntingan terakhir sebelum naskah diajukan.

\begin{ilustrasi}
Rina menyusun Laporan Kegiatan dari rekaman sebuah kegiatan lapangan. Ia mengisi kepala dengan tanggal dan tempat yang ia ketahui, lalu mengosongkan kolom disampaikan kepada karena belum diputuskan. Naskah pertamanya ditolak pemeriksa karena salah satu angka pada bagian hasil tidak memiliki keterangan siapa yang menyebutnya. Rina kembali ke transkrip, menemukan penyebut angka itu, lalu menyusun ulang naskahnya. Naskah kedua lolos, dan ia menyunting berkas Word-nya sebelum diajukan kepada pimpinan.
\end{ilustrasi}

\begin{kegiatan}{Latihan 5. Menyusun dokumen dan membaca penolakan (20 menit)}
Dari transkrip latihan, susun satu Laporan Kegiatan. Isi kepala hanya dengan informasi yang Anda temukan di transkrip. Setelah selesai, periksa bagian III: pastikan setiap angka disertai penyebutnya. Bila naskah ditolak, tulis alasannya dan rumuskan satu langkah perbaikan. Keluaran yang diharapkan: berkas \texttt{.docx} hasil Anda dan satu kalimat tentang perbaikan yang dilakukan.
\end{kegiatan}

\begin{rangkuman}
Dokumen resmi disusun dari dua jenis naskah dengan bagian yang tetap. Kepala dokumen diisi pengguna, dan Pandora tidak mengarang tanggal, tempat, atau nama. Setiap angka harus disertai penyebut, dan bagian yang tidak jelas ditulis sebagai \texttt{[tidak jelas]}. Naskah yang ditolak pemeriksa perlu disusun ulang sesuai alasannya. Rina memperbaiki naskah pertamanya, dan naskah kedua lolos sebelum diajukan.
\end{rangkuman}

\chapter{Riwayat dan Masalah Umum}

Setiap pekerjaan di Pandora tersimpan sebagai folder hasil di riwayat akun Anda. Riwayat memudahkan Anda membuka kembali hasil lama, melanjutkan pekerjaan yang belum selesai, dan menghapus hasil yang tidak diperlukan. Bab ini membahas cara mengelola riwayat dan cara menangani masalah yang paling sering muncul. Masalah yang tidak tercantum di sini dilaporkan kepada pengelola bersama langkah yang sudah Anda coba.

\begin{indikator}
Setelah mempelajari bab ini, peserta mampu:
\begin{enumerate}
  \item membuka kembali hasil lama dari riwayat dan melanjutkan pekerjaan yang belum selesai;
  \item menghapus hasil dengan aman setelah memahami akibatnya;
  \item mengenali lima masalah umum dan mengambil langkah pertama untuk masing-masing.
\end{enumerate}
\end{indikator}

\section{Membuka dan Melanjutkan Riwayat}
Panel riwayat menampilkan semua folder hasil milik Anda, dengan nama yang mengikuti pola \texttt{01\_nama}, \texttt{02\_nama}, dan seterusnya. Klik sebuah folder untuk membuka hasilnya kembali, termasuk transkrip, rapi AI, notulensi, dan dokumen bila sudah dibuat.

Bila sebuah folder belum memiliki notulensi atau dokumen, tombol untuk membuatnya tetap tersedia setelah folder dibuka. Dengan cara ini, Anda tidak perlu mentranskripsi ulang rekaman yang sama.

\section{Menghapus Riwayat}
Menghapus folder hasil berarti menghapus semua berkas di dalamnya, yaitu transkrip, hasil rapi, notulensi, dan dokumen. Pandora meminta konfirmasi sebelum menghapus. Pesan konfirmasi menyebut nama folder yang akan dihapus. Penghapusan tidak dapat dibatalkan.

Sebelum menghapus, pastikan Anda tidak lagi memerlukan hasilnya, dan bila perlu unduh dokumen yang masih dibutuhkan. Jangan menghapus hasil yang masih Anda perlukan sebagai rujukan, misalnya transkrip asli yang menjadi sumber notulensi atau dokumen yang sudah diajukan.

\section{Masalah Umum dan Langkah Pertama}
Tabel~\ref{tab:masalah} merangkum masalah yang paling sering muncul beserta langkah pertama yang bisa Anda lakukan.

\begin{table}[H]
\caption{Masalah umum dan langkah pertama}
\label{tab:masalah}
{\small
\begin{tabularx}{\linewidth}{@{}P{4.2cm}YY@{}}
\toprule
\rowcolor{ugmpanel}\thd{Gejala} & \thd{Kemungkinan penyebab} & \thd{Langkah pertama} \\
\midrule
Tombol AI tidak dapat diklik & Layanan AI belum siap & Periksa status di bagian atas layar, lalu hubungi pengelola \\
Pesan worker sedang digunakan & Transkripsi lain sedang berjalan & Tunggu sampai selesai, lalu coba lagi \\
Berkas tidak dapat diunggah & Format berkas tidak didukung & Ubah ke format yang didukung, lalu unggah ulang \\
Log berulang atau kosong di bagian tertentu & Bagian hening atau suara tidak jelas & Periksa rekaman di menit tersebut sebelum menyimpulkan \\
Notulensi atau dokumen tidak selesai & Layanan AI sibuk atau terhenti & Tunggu, lalu coba lagi. Bila tetap gagal, catat pesan dan hubungi pengelola \\
Naskah ditolak pemeriksa & Cakupan, bentuk, atau keterlacakan tidak terpenuhi & Baca alasan yang ditampilkan, lalu susun ulang sesuai bagian \\
\bottomrule
\end{tabularx}
}
\end{table}

Bila masalah masih berlanjut setelah langkah pertama, catat tiga hal: waktu kejadian, pesan yang muncul, dan nama folder hasil yang terkait. Informasi ini membantu pengelola menelusuri masalah tanpa perlu menanyakan ulang.

\begin{ilustrasi}
Pada hari kerja yang padat, Rina mencoba memulai transkripsi rapat kedua, tetapi layar menampilkan pesan bahwa worker sedang digunakan. Ia tidak mengulang berkali-kali. Ia menunggu, lalu mencoba lagi sekitar dua puluh menit kemudian. Setelah itu, ia membuka riwayat dan menghapus folder percobaan pertamanya yang tidak lagi diperlukan, dengan membaca pesan konfirmasi terlebih dahulu.
\end{ilustrasi}

\begin{kegiatan}{Latihan 6. Menelusuri riwayat dan masalah (10 menit)}
Buka riwayat dan pilih satu folder hasil dari latihan sebelumnya. Tentukan apakah notulensi dan dokumen sudah ada di dalamnya. Lalu pilih satu gejala dari Tabel~\ref{tab:masalah} dan tulis langkah pertama yang akan Anda ambil, tanpa benar-benar menghapus apa pun. Keluaran yang diharapkan: daftar kondisi folder hasil dan satu langkah pertama yang Anda pilih.
\end{kegiatan}

\begin{rangkuman}
Riwayat menyimpan semua folder hasil milik akun Anda dan memungkinkan Anda melanjutkan pekerjaan lama. Menghapus folder berarti menghapus seluruh isinya, sehingga Pandora meminta konfirmasi terlebih dahulu. Sebagian besar masalah umum dapat diatasi dengan menunggu, memeriksa status, atau menyusun ulang naskah sesuai alasannya. Rina menunggu giliran transkripsi dan menghapus hasil yang tidak lagi diperlukan. Bab berikutnya menutup modul dengan kembali ke pertanyaan pemantik.
\end{rangkuman}

\chapter{Penutup}

\section{Kerangka Berpikir}
Seluruh modul ini mengikuti satu urutan. Rekaman menjadi transkrip asli, dan transkrip asli menjadi sumber bagi semua hasil lain. Setiap turunan, baik hasil rapi, notulensi, maupun dokumen resmi, harus dicocokkan kembali dengan sumbernya sebelum dibagikan. Pandora membantu mempercepat pekerjaan, tetapi tanggung jawab atas isi tetap berada pada pengguna dan pejabat yang berwenang.

\section{Kembali ke Pertanyaan Pemantik}
Pada Bab I, Anda telah menuliskan jawaban awal atas pertanyaan tentang rapat atau kegiatan yang paling sulit dicatat ulang, siapa yang paling membutuhkan catatannya, dan bagaimana Anda tahu catatan itu benar. Bandingkan jawaban itu dengan langkah-langkah yang sudah Anda pelajari, lalu catat apa yang kini terpikirkan tetapi belum muncul pada awal pelatihan.

\section{Daftar Periksa Sebelum Membagikan Hasil}
\begin{ilustrasi}
Tiga minggu setelah pelatihan, Rina diminta menyusun laporan kegiatan lapangan dan notulensi rapat kerja untuk disampaikan kepada pimpinan unit. Ia membuka Pandora dengan langkah yang sudah ia kenal. Sebelum mengirim, ia memakai daftar periksa di bawah ini. Pada kali pertama, ia menemukan satu keputusan yang tercatat sebagai usulan dan memperbaikinya. Ia tidak lagi merasa harus mendengarkan seluruh rekaman dari awal, tetapi ia tetap membuka menit-menit kunci.
\end{ilustrasi}

Periksa setiap hal berikut sebelum hasil dibagikan kepada pihak lain.

\begin{enumerate}
  \item \textbf{Sumber.} Setiap angka, nama, dan tanggal dapat ditelusuri ke transkrip asli.
  \item \textbf{Keputusan.} Hanya hal yang benar-benar diputuskan yang ditulis sebagai keputusan.
  \item \textbf{Penyebut.} Setiap angka disertai keterangan siapa yang menyebutnya.
  \item \textbf{Bagian tidak jelas.} Tidak ada tebakan. Bagian yang tidak terdengar ditandai.
  \item \textbf{Kepala dan identitas.} Kolom kepala hanya berisi informasi yang Anda ketahui.
  \item \textbf{Pengesahan.} Dokumen yang memerlukan pengesahan tidak dibagikan sebagai dokumen final sebelum disahkan.
  \item \textbf{Keterbukaan.} Bila hasil dibuat dengan bantuan AI, catat keterangan itu sesuai ketentuan instansi.
\end{enumerate}

\section{Tindak Lanjut di Unit Kerja}
Setelah pelatihan, ada empat langkah yang dapat langsung Anda kerjakan di unit kerja masing-masing:
\begin{enumerate}
  \item menentukan satu rapat atau kegiatan pertama yang akan dicatat dengan Pandora, lalu menyepakati siapa yang memeriksa hasilnya;
  \item menyusun aturan internal sederhana tentang rekaman apa saja yang boleh diunggah, sesuai Ketentuan Penggunaan Data;
  \item membuat folder kerja bersama untuk berkas Word yang sudah diperiksa, bukan untuk rekaman mentah;
  \item menyampaikan kepada pengelola masalah yang belum tercantum dalam modul ini, agar modul dapat diperbarui.
\end{enumerate}

\begin{rangkuman}
Pandora mengubah rekaman menjadi transkrip, lalu membantu menyusun hasil turunan yang rapi dan dapat dibagikan. Transkrip asli menjadi rujukan utama, dan setiap turunan perlu diperiksa terhadap sumbernya. Notulensi dan dokumen resmi memiliki aturan yang ketat, terutama untuk keputusan, angka, nama, dan kepala dokumen. Rina memulai dari rekaman yang tidak ia kenal, dan pada akhirnya ia mampu menyusun, memeriksa, dan membagikan hasilnya sendiri dengan percaya diri.
\end{rangkuman}

\chapter*{Daftar Pustaka}
\addcontentsline{toc}{chapter}{Daftar Pustaka}

% Gaya APA ringkas, urut abjad. Peraturan ditulis lengkap tanpa tahun di depan.
\begin{itemize}[leftmargin=1cm, itemindent=-1cm, label={}, itemsep=8pt]
  \item Kementerian Komunikasi dan Informatika Republik Indonesia. (2023). Surat Edaran Menteri Komunikasi dan Informatika Nomor 9 Tahun 2023 tentang Etika Kecerdasan Artifisial.
  \item Presiden Republik Indonesia. (2018). \textit{Peraturan Presiden Nomor 95 Tahun 2018 tentang Sistem Pemerintahan Berbasis Elektronik}.
  \item Radford, A., Kim, J. W., Xu, T., Brockman, G., McLeavey, C., dan Sutskever, I. (2023). Robust speech recognition via large-scale weak supervision. Dalam \textit{Proceedings of the 40th International Conference on Machine Learning} (PMLR 202).
  \item Tim Pengelola Pandora. (2026). \textit{Panduan penggunaan Pandora} [Dokumen internal].
\end{itemize}


\end{document}
